% ============================================================================
% Fichier  : partie3/P03_C11_droit.tex
% Titre    : Le Droit de la Sécurité au Cameroun et en Afrique
% Auteur   : MINKA MI NGUIDJOÏ Thierry Emmanuel
% Version  : 1.0 -- novembre 2026
% ============================================================================
% Loi n° 2010/012 lue dans le texte officiel : art. 7 (1), 25 (1), 42,
% 65 (2), 68, 69 cités. Loi n° 2010/013 et loi n° 2024/017 présentées sans
% citation d'article (à collationner). Conventions de Budapest (2001) et de
% Malabo (2014) : les ratifications par le Cameroun ne sont pas affirmées ;
% elles sont à vérifier sur les sites des dépositaires.
% ============================================================================

\chapter{Le Droit de la Sécurité au Cameroun et en Afrique}
\label{chap:droit}
\chaptermark{Droit de la sécurité}

\epigraph{%
    « Le droit ne protège pas à la place des mesures~;\\
    il dit qui répond quand elles manquent. »%
}{Principe du cours}

\niveauChapitre{Fondamental}{Chapitres~\ref{chap:proprietes}
et~\ref{chap:gouvernance}.}

\begin{boxobjectifs}
\begin{itemize}
    \item Présenter les textes camerounais de la sécurité numérique et le
          rôle de l'ANTIC.
    \item Identifier les obligations des exploitants (sécurité,
          conservation, confidentialité).
    \item Qualifier les principales infractions informatiques de la loi
          n°~2010/012.
    \item Situer les conventions africaine et internationale.
    \item Savoir quoi faire, juridiquement, après une attaque.
\end{itemize}
\end{boxobjectifs}

\bigskip

% ============================================================================
\section*{Question directrice : que dit la loi sur ce qui s'est passé à Garoua~?}
% ============================================================================

Le lundi, le directeur juridique de LiZbiZ pose quatre questions~:
l'attaquant a-t-il commis une infraction, et laquelle~? LiZbiZ a-t-elle
manqué à une obligation~? Doit-elle conserver ou remettre des journaux~?
Peut-elle payer la rançon~? Ce chapitre donne les textes pour y répondre,
en citant seulement ce qui a été lu dans sa version officielle.

\medskip
\textbf{Plan du chapitre.}
\begin{enumerate}
    \item Les textes et les institutions (section~\ref{sec:c11-textes}).
    \item Les obligations des exploitants (section~\ref{sec:c11-obligations}).
    \item Les infractions (section~\ref{sec:c11-infractions}).
    \item Les données personnelles (section~\ref{sec:c11-donnees}).
    \item Le cadre africain et international (section~\ref{sec:c11-international}).
    \item Après une attaque (section~\ref{sec:c11-apres}).
\end{enumerate}

\separateur

% ============================================================================
\section{Les Textes et les Institutions}
\label{sec:c11-textes}
% ============================================================================

\begin{tableaucours}{L{4.4cm}L{7.8cm}}{%
    \tableauHeader{Texte} &
    \tableauHeader{Objet}
}
Loi n°~2010/012 du 21~décembre~2010 & Cybersécurité et cybercriminalité~:
définitions, obligations des exploitants, certification, infractions,
procédure \\
Loi n°~2010/013 du 21~décembre~2010 & Communications électroniques~;
institue l'ANTIC \\
Loi n°~2024/017 du 23~décembre~2024 & Protection des données à caractère
personnel (articles à collationner au Journal officiel) \\
\end{tableaucours}
\vspace{-0.8em}
\begin{center}{\small\itshape Tableau~11.1. Les textes camerounais}\end{center}

\begin{boxjuris}
\textbf{\loicam{n°~2010/012}, article~7, alinéa~1} (texte officiel)~:
\textit{« L'Agence Nationale des Technologies de l'Information et de la
Communication, ci-après désignée l'Agence, instituée par la loi régissant
les communications électroniques au Cameroun, est chargée de la
régulation des activités de sécurité électronique, en collaboration avec
l'Agence de Régulation des Télécommunications. »}

\medskip
\textbf{Analyse.} L'ANTIC régule, contrôle et suit la sécurité des
systèmes d'information~; elle est aussi l'autorité de certification racine
(art.~8, Chapitre~\ref{chap:crypto-asym}) et conduit les audits de
sécurité obligatoires des opérateurs (art.~13 et~32).
\end{boxjuris}

% ============================================================================
\section{Les Obligations des Exploitants}
\label{sec:c11-obligations}
% ============================================================================

\begin{boxjuris}
\textbf{\loicam{n°~2010/012}} (texte officiel)

\medskip
\textbf{Article~25, alinéa~1}~: \textit{« Les opérateurs de réseaux et
les fournisseurs de services de communications électroniques ont
l'obligation de conserver les données de connexion et de trafic pendant
une période de dix (10) ans. »}

\medskip
\textbf{Article~42}~: \textit{« La confidentialité des communications
acheminées à travers les réseaux de communications électroniques et les
systèmes d'information, y compris les données relatives au trafic, est
assurée par les opérateurs et exploitants des réseaux de communications
électroniques et des systèmes d'information. »}

\medskip
\textbf{Analyse.} À ces obligations s'ajoute celle de l'article~26
(gestion continue et normalisée des risques, Chapitre~\ref{chap:gouvernance}),
et l'article~29 impose aux exploitants de systèmes d'information de
conserver les données de connexion et de trafic de leurs systèmes. La
conservation des journaux est donc une obligation légale, pas seulement
une bonne pratique~: une entreprise sans journaux ne peut ni enquêter ni
prouver qu'elle a respecté la loi.
\end{boxjuris}

% ============================================================================
\section{Les Infractions}
\label{sec:c11-infractions}
% ============================================================================

\begin{boxjuris}
\textbf{\loicam{n°~2010/012}} (texte officiel)

\medskip
\textbf{Article~65, alinéa~2}~: \textit{« Est puni des peines prévues à
l'alinéa 1 ci-dessus, tout accès non autorisé, à l'ensemble ou à une
partie d'un réseau de communications électroniques ou d'un système
d'information ou d'un équipement terminal. »}

\medskip
\textbf{Article~68, alinéa~1}~: \textit{« Est puni d'un emprisonnement de
cinq (05) à dix (10) ans et d'une amende de 10.000.000 (dix millions) à
50.000.000 (cinquante millions) F CFA ou de l'une de ces deux peines
seulement, celui qui accède ou se maintient, frauduleusement, dans tout ou
partie d'un réseau de communications électroniques ou d'un système
d'information en transmettant, endommageant, provoquant une perturbation
grave ou une interruption du fonctionnement dudit système ou dudit
réseau. »} \textbf{Alinéa~2}~: \textit{« Les peines prévues à l'alinéa 1
ci-dessus sont doublées s'il en est résulté, soit la suppression ou la
modification des données contenues dans le système d'information, soit
une altération de son fonctionnement. »}

\medskip
\textbf{Article~69}~: \textit{« Est puni d'un emprisonnement de cinq (05)
à dix (10) ans et d'une amende de 10.000.000 (dix millions) à
100.000.000 (cent millions) F CFA ou de l'une de ces deux peines
seulement, celui qui accède sans droit, et en violation des mesures de
sécurité, à l'ensemble ou à une partie d'un réseau de communications
électroniques, d'un système d'information ou d'un équipement terminal,
afin d'obtenir des informations ou des données, en relation avec un
système d'information connecté à un autre système d'information. »}
\end{boxjuris}

\begin{boxscenario}{Opération GAROUA. Quelles infractions~?}
\textbf{Faits.} Accès avec un mot de passe volé, maintien pendant trois
jours, copie de données, suppression des sauvegardes, chiffrement des
serveurs.

\medskip
\textbf{Analyse} (confiance moyenne). Les faits paraissent relever de
l'accès et du maintien frauduleux avec perturbation et suppression de
données (art.~68, peines doublées par l'alinéa~2), et de l'accès en
violation des mesures de sécurité pour obtenir des données (art.~69).
\textbf{La qualification appartient au juge}~: le RSSI décrit les faits et
conserve les preuves, il ne qualifie pas (Chapitre~\ref{chap:investigation}).
\end{boxscenario}

\begin{boxalert}
\textbf{Violation typique.} Un administrateur « teste » le système d'un
partenaire, sans autorisation écrite, pour lui prouver qu'il est
vulnérable.

\smallskip
\textbf{Règle~:} l'accès non autorisé est puni (art.~65, al.~2), quelle
que soit l'intention. Tout test se fait avec une autorisation écrite du
responsable du système (Chapitres~\ref{chap:ethique} et~\ref{chap:evaluer}).
\end{boxalert}

% ============================================================================
\section{Les Données Personnelles}
\label{sec:c11-donnees}
% ============================================================================

La loi n°~2010/012 pose déjà le droit au respect de la vie privée (art.~41)
et réprime l'atteinte à l'intimité de la vie privée par la captation ou la
transmission de données privées sans consentement (art.~74). La loi
n°~2024/017 organise la protection des données à caractère personnel~:
principes de traitement (finalité, proportionnalité, durée de
conservation, sécurité), droits des personnes, obligations des
responsables de traitement. Ses articles, notamment sur la notification
des violations de données, sont à lire dans le texte officiel avant toute
mise en œuvre~; le cours ne les cite pas sans les avoir collationnés.

% ============================================================================
\section{Le Cadre Africain et International}
\label{sec:c11-international}
% ============================================================================

\begin{tableaucours}{L{4.2cm}L{8.0cm}}{%
    \tableauHeader{Texte} &
    \tableauHeader{Apport}
}
Convention de Malabo (Union africaine, 2014) & Cybersécurité et protection
des données à caractère personnel~; harmonisation des incriminations sur le
continent~\cite{malabo2014} \\
Convention de Budapest (Conseil de l'Europe, 2001) & Premier traité sur la
cybercriminalité~; incriminations, procédure, entraide
internationale~\cite{budapest2001} \\
\end{tableaucours}
\vspace{-0.8em}
\begin{center}{\small\itshape Tableau~11.2. Les conventions internationales}\end{center}

\begin{boxinfo}
\textbf{Vérifier l'état des ratifications.} L'application d'une
convention dépend de la ratification par l'État. Avant de l'invoquer,
on vérifie l'état des signatures et des ratifications sur le site du
dépositaire (Union africaine, Conseil de l'Europe).
\end{boxinfo}

% ============================================================================
\section{Après une Attaque}
\label{sec:c11-apres}
% ============================================================================

\begin{boxprocedure}
\textbf{Les réflexes juridiques après un incident}
\begin{enumerate}
    \item \textbf{Conserver les preuves}~: journaux, images des systèmes,
          messages, avec leurs empreintes (Chapitre~\ref{chap:investigation}).
    \item \textbf{Porter plainte} auprès des services de police ou de
          gendarmerie compétents, ou du procureur, en remettant une
          description factuelle.
    \item \textbf{Informer} les autorités et les personnes dans les cas où
          un texte l'exige, en vérifiant les délais dans le texte officiel.
    \item \textbf{Associer le conseil juridique} à toute décision sur la
          rançon et sur la communication.
    \item \textbf{Ne pas qualifier} publiquement les faits ni désigner un
          coupable.
\end{enumerate}
\end{boxprocedure}

\begin{boxinfo}
\textbf{Payer la rançon~?} Le paiement n'assure ni la restitution des
données, ni l'absence de publication, et il finance l'activité
criminelle. Les autorités de nombreux pays le déconseillent. La décision,
si elle se pose, appartient à la direction, avec son conseil juridique,
et se documente.
\end{boxinfo}

\begin{boxmetier}{RSSI}
\textbf{Connaître ses obligations.} Le RSSI tient la liste des
obligations légales applicables (conservation, confidentialité,
notification, audit), avec pour chacune la mesure qui la satisfait et la
preuve.
\end{boxmetier}

\begin{boxmetier}{Administrateur systèmes et réseaux}
\textbf{Rester dans son droit.} L'administrateur n'accède qu'aux systèmes
dont il a la charge, ne lit pas les communications qu'il achemine
(art.~42 et~44), et ne teste un système tiers qu'avec une autorisation
écrite.
\end{boxmetier}

\separateur

\begin{boxresume}
\textbf{Ce qu'il faut retenir de ce chapitre~:}
\begin{itemize}
    \item \textbf{La loi n°~2010/012} définit, oblige et punit~; l'ANTIC
          régule la sécurité électronique (art.~7).
    \item \textbf{Les exploitants} conservent les données de connexion et de
          trafic (art.~25 et~29) et assurent la confidentialité des
          communications (art.~42).
    \item \textbf{L'accès non autorisé} est puni (art.~65), plus sévèrement
          avec perturbation ou suppression de données (art.~68) ou pour
          obtenir des données (art.~69).
    \item \textbf{La loi n°~2024/017} encadre les données personnelles~; ses
          articles se lisent dans le texte officiel.
    \item \textbf{Malabo et Budapest} structurent la coopération~; on
          vérifie les ratifications.
    \item \textbf{Après une attaque}~: conserver, porter plainte, informer
          si la loi l'exige, ne pas qualifier.
\end{itemize}
\end{boxresume}

\begin{boxinfo}
\textbf{Points de contrôle de l'élève.} Je sais~: citer le rôle de
l'ANTIC~; énoncer les obligations des articles 25 et~42~; distinguer les
articles 65, 68 et~69~; dire ce que je dois vérifier avant d'appliquer la
loi n°~2024/017 ou une convention~; énumérer les réflexes juridiques après
un incident.
\end{boxinfo}

\bigskip

\begin{boxexo}[title={\ding{45}\quad Exercice~11.1. Les institutions \niveaubase}]
Quel est le rôle de l'ANTIC selon les articles 7 et~8 de la loi
n°~2010/012~?
\end{boxexo}

\begin{boxexo}[title={\ding{45}\quad Exercice~11.2. Conserver \niveaubase}]
Quelle durée de conservation des données de connexion la loi impose-t-elle
aux opérateurs~? Pourquoi cette obligation sert-elle l'enquête~?
\end{boxexo}

\begin{boxexo}[title={\ding{45}\quad Exercice~11.3. Qualifier \niveaumoyen}]
Pour chaque fait, indiquez l'article qui paraît applicable, en rappelant
que la qualification appartient au juge~: (a) un étudiant entre dans le
serveur de son école avec un mot de passe trouvé~; (b) un ancien salarié
se connecte et efface la base des clients~; (c) un concurrent pénètre le
système pour copier des tarifs.
\end{boxexo}

\begin{boxexo}[title={\ding{45}\quad Exercice~11.4. Doubler la peine \niveaumoyen}]
Dans quel cas l'article~68, alinéa~2, double-t-il les peines~? Appliquez-le
à l'Opération GAROUA.
\end{boxexo}

\begin{boxexo}[title={\ding{45}\quad Exercice~11.5. Le test non autorisé \niveaumoyen}]
Un étudiant balaye les ports du réseau de l'ENSPY « pour son mémoire »,
sans autorisation. Analysez la situation au regard de l'article~65.
\end{boxexo}

\begin{boxexo}[title={\ding{45}\quad Exercice~11.6. Confidentialité \niveaumoyen}]
Un administrateur de la BANQUE SANAGA lit les courriels d'un collègue
« pour vérifier une rumeur ». Quels articles sont en cause~?
\end{boxexo}

\begin{boxexo}[title={\ding{45}\quad Exercice~11.7. Avant d'appliquer \niveaumoyen}]
Votre directeur vous demande « ce que dit la loi 2024/017 sur la
notification des fuites ». Que faites-vous avant de répondre, et pourquoi~?
\end{boxexo}

\begin{boxexo}[title={\ding{45}\quad Exercice~11.8. La rançon \niveaumoyen}]
Préparez, pour le comité de LiZbiZ, les arguments pour et contre le
paiement de la rançon, et les questions juridiques à poser au conseil.
\end{boxexo}

\begin{boxexo}[title={\ding{45}\quad Exercice~11.9. Le dossier de plainte \niveauexpert}]
Rédigez la description factuelle des faits de l'Opération GAROUA à
joindre à la plainte~: sans qualification, avec les preuves disponibles.
\end{boxexo}

\begin{boxexo}[title={\ding{45}\quad Exercice~11.10. Le registre des obligations \niveauexpert}]
Construisez le registre des obligations légales de sécurité de LiZbiZ~:
obligation, texte, mesure, preuve, responsable.
\end{boxexo}

\bigskip

\begin{boxinfo}
\textbf{Pour aller plus loin.}
\begin{itemize}
    \item Loi n°~2010/012 relative à la cybersécurité et à la
          cybercriminalité~\cite{loi2010012}.
    \item Convention de Malabo~\cite{malabo2014}~; Convention de
          Budapest~\cite{budapest2001}.
\end{itemize}
\end{boxinfo}

\begin{boxinfo}
\textbf{Transition.} Le droit trace des limites~; l'éthique demande de
bien agir à l'intérieur. Le Chapitre~\ref{chap:ethique} aborde la
responsabilité du professionnel de la sécurité.
\end{boxinfo}
